\section{Verification index}
\label{app:verification}

Table~\ref{tab:results} lists every registered analytic result with its
recorded status and evidence level. The two columns are independent
fields and are meant to be read against each other: \texttt{proven} with
\texttt{record-backed} evidence means the argument exists and the machine
artefact does not.

Of \ResTotal{} registered results, \ResProven{} are recorded
\texttt{proven}, \ResConditional{} \texttt{conditional},
\ResDisputed{} \texttt{disputed} and \ResFalsified{} \texttt{falsified}.
The Lean tree carries \LeanTotal{} registered theorems with
\LeanSorry{} occurrences of \texttt{sorry}.

A caution on both counts. \ResProven{} out of \ResTotal{} is a high
proportion because a claim is not registered here until its argument is
written down; the register measures what has been recorded, not what
fraction of the problem is solved. Likewise \LeanTotal{} theorems is a
count of declarations, not of mathematical content --- see
\S\ref{app:repro}, which explains why a theorem count is not a coverage
measure and cites a case in this program's own archive where seventy-one
\texttt{sorry}-free modules turned out to prove tautologies.

% Generated by paper/master/generate_apparatus.py.  Do not edit.
\begin{longtable}{@{}p{0.44\linewidth} l l@{}}
\caption{Analytic results by recorded status and evidence level. \emph{Status} is what the argument establishes; \emph{evidence} is what artefact backs it. The two are deliberately independent: a result can be \texttt{proven} in status and \texttt{record-backed} in evidence, meaning the proof exists and the machine artefact does not.}\label{tab:results}\\
\toprule
Result & Status & Evidence \\
\midrule
\endfirsthead
\toprule Result & Status & Evidence \\ \midrule \endhead
\bottomrule
\endfoot
\texttt{\scriptsize ANISOTROPY\_\allowbreak{}THREE\_\allowbreak{}COORDINATE\_\allowbreak{}SHARP\_\allowbreak{}BOUND} & proven & analytic \\
\texttt{\scriptsize ARCHIVE\_\allowbreak{}CENTER\_\allowbreak{}DEFECT\_\allowbreak{}OBSTRUCTIONS} & proven & analytic \\
\texttt{\scriptsize ARCHIVE\_\allowbreak{}CONCENTRATION\_\allowbreak{}LOCALIZATION} & proven & analytic \\
\texttt{\scriptsize ARCHIVE\_\allowbreak{}FINITE\_\allowbreak{}RANGE\_\allowbreak{}INVERSE} & proven & analytic \\
\texttt{\scriptsize ARCHIVE\_\allowbreak{}GLOBAL\_\allowbreak{}PAIRING\_\allowbreak{}OBSTRUCTION} & proven & analytic \\
\texttt{\scriptsize ARCHIVE\_\allowbreak{}HORIZONTAL\_\allowbreak{}RESTRICTION\_\allowbreak{}BOUNDARY} & proven & analytic \\
\texttt{\scriptsize ARCHIVE\_\allowbreak{}LYAPUNOV\_\allowbreak{}PATCHING} & proven & analytic \\
\texttt{\scriptsize ARCHIVE\_\allowbreak{}REFLECTION\_\allowbreak{}PERMANENCE} & proven & analytic \\
\texttt{\scriptsize ARCHIVE\_\allowbreak{}SIGNED\_\allowbreak{}SU2\_\allowbreak{}DIFFUSION} & proven & analytic \\
\texttt{\scriptsize ARCHIVE\_\allowbreak{}SMOOTH\_\allowbreak{}PROFILES} & proven & analytic \\
\texttt{\scriptsize ARCHIVE\_\allowbreak{}SOURCE\_\allowbreak{}TILT\_\allowbreak{}AND\_\allowbreak{}ROOTED\_\allowbreak{}BOUNDS} & proven & analytic \\
\texttt{\scriptsize ASSEMBLED\_\allowbreak{}Q4\_\allowbreak{}COERCIVITY} & proven & analytic \\
\texttt{\scriptsize BALABAN\_\allowbreak{}MULTISCALE\_\allowbreak{}CAUCHY\_\allowbreak{}SUMMABILITY} & proven & analytic \\
\texttt{\scriptsize COMMON\_\allowbreak{}SPACE\_\allowbreak{}CONTRACTIVE\_\allowbreak{}DRIFT} & proven & analytic \\
\texttt{\scriptsize COMPACT\_\allowbreak{}ROTOR\_\allowbreak{}FAST\_\allowbreak{}GAP} & proven & analytic \\
\texttt{\scriptsize CONDITIONAL\_\allowbreak{}GRADIENT\_\allowbreak{}REPAIR} & proven & analytic \\
\texttt{\scriptsize CONDITIONAL\_\allowbreak{}QUANTUM\_\allowbreak{}PATH\_\allowbreak{}COVARIANCE} & proven & analytic \\
\texttt{\scriptsize CREATOR\_\allowbreak{}PARENT\_\allowbreak{}GAP} & proven & analytic \\
\texttt{\scriptsize CREATOR\_\allowbreak{}VELOCITY\_\allowbreak{}INVERSION} & proven & analytic \\
\texttt{\scriptsize DIMENSION\_\allowbreak{}FIVE\_\allowbreak{}OPERATOR\_\allowbreak{}IRRELEVANCE} & proven & analytic \\
\texttt{\scriptsize DYNAMIC\_\allowbreak{}CONDITIONAL\_\allowbreak{}CUBIC\_\allowbreak{}ENERGY} & proven & analytic \\
\texttt{\scriptsize EARLIER\_\allowbreak{}B6\_\allowbreak{}RETAINED\_\allowbreak{}KRYLOV\_\allowbreak{}CAMPAIGN} & conditional & numerical \\
\texttt{\scriptsize EARLIER\_\allowbreak{}CONDITIONAL\_\allowbreak{}SPECTRAL\_\allowbreak{}FLOOR} & proven & analytic \\
\texttt{\scriptsize EARLIER\_\allowbreak{}CORNER\_\allowbreak{}BLOCKING\_\allowbreak{}REFLECTION\_\allowbreak{}DEFECT} & proven & analytic \\
\texttt{\scriptsize EARLIER\_\allowbreak{}CUBIC\_\allowbreak{}QUOTIENT\_\allowbreak{}REGULARITY} & proven & analytic \\
\texttt{\scriptsize EARLIER\_\allowbreak{}D4\_\allowbreak{}DISJOINT\_\allowbreak{}STAPLE\_\allowbreak{}COORDINATES} & proven & analytic \\
\texttt{\scriptsize EARLIER\_\allowbreak{}DECORATED\_\allowbreak{}HISTORY\_\allowbreak{}MERGING} & proven & analytic \\
\texttt{\scriptsize EARLIER\_\allowbreak{}DETERMINANT\_\allowbreak{}PARITY\_\allowbreak{}REPAIR} & proven & analytic \\
\texttt{\scriptsize EARLIER\_\allowbreak{}EQUAL\_\allowbreak{}RAY\_\allowbreak{}IDENTIFICATION} & proven & analytic \\
\texttt{\scriptsize EARLIER\_\allowbreak{}FINITE\_\allowbreak{}ORDER\_\allowbreak{}SUPPORT} & proven & analytic \\
\texttt{\scriptsize EARLIER\_\allowbreak{}GRAM\_\allowbreak{}NULL\_\allowbreak{}OPERATOR\_\allowbreak{}QUOTIENT} & proven & analytic \\
\texttt{\scriptsize EARLIER\_\allowbreak{}ISOLATED\_\allowbreak{}POSITIVE\_\allowbreak{}MATRIX\_\allowbreak{}ATOM} & proven & analytic \\
\texttt{\scriptsize EARLIER\_\allowbreak{}JOINT\_\allowbreak{}RANK\_\allowbreak{}VOLUME\_\allowbreak{}SCALING} & proven & analytic \\
\texttt{\scriptsize EARLIER\_\allowbreak{}LOCALIZED\_\allowbreak{}FRAME\_\allowbreak{}OBSTRUCTION} & proven & analytic \\
\texttt{\scriptsize EARLIER\_\allowbreak{}MOVING\_\allowbreak{}ADJOINT\_\allowbreak{}FORCE\_\allowbreak{}DERIVATIVE} & proven & analytic \\
\texttt{\scriptsize EARLIER\_\allowbreak{}NESTED\_\allowbreak{}QUOTIENT\_\allowbreak{}REDUCTION} & proven & analytic \\
\texttt{\scriptsize EARLIER\_\allowbreak{}POSITIVE\_\allowbreak{}SECTOR\_\allowbreak{}PHASE\_\allowbreak{}ISOLATION} & proven & analytic \\
\texttt{\scriptsize EARLIER\_\allowbreak{}REFLECTION\_\allowbreak{}ADAPTED\_\allowbreak{}BLOCKING} & proven & analytic \\
\texttt{\scriptsize EARLIER\_\allowbreak{}SHELL\_\allowbreak{}SYMMETRY\_\allowbreak{}RITZ\_\allowbreak{}CONTROL} & proven & analytic \\
\texttt{\scriptsize EARLIER\_\allowbreak{}SPHERE\_\allowbreak{}EQUAL\_\allowbreak{}COUPLING\_\allowbreak{}COVARIANCE} & proven & analytic \\
\texttt{\scriptsize EARLIER\_\allowbreak{}SU3\_\allowbreak{}POSITIVITY\_\allowbreak{}BOUNDARIES} & proven & analytic \\
\texttt{\scriptsize EARLIER\_\allowbreak{}VSU\_\allowbreak{}BOUNDED\_\allowbreak{}DOMAIN\_\allowbreak{}WELLPOSEDNESS} & proven & analytic \\
\texttt{\scriptsize EXPONENTIAL\_\allowbreak{}SOURCE\_\allowbreak{}TOTALITY} & proven & analytic \\
\texttt{\scriptsize FESHBACH\_\allowbreak{}RESOLVENT\_\allowbreak{}COMPARISON} & proven & analytic \\
\texttt{\scriptsize FORM\_\allowbreak{}SCHUR\_\allowbreak{}SCALE\_\allowbreak{}COMPARISON} & proven & analytic \\
\texttt{\scriptsize FULL\_\allowbreak{}GRAPH\_\allowbreak{}SOURCE\_\allowbreak{}CONGRUENCE} & proven & analytic \\
\texttt{\scriptsize G17\_\allowbreak{}CORRECTED\_\allowbreak{}KP\_\allowbreak{}POLYMER\_\allowbreak{}STABILITY} & proven & analytic \\
\texttt{\scriptsize G9\_\allowbreak{}COMBINED\_\allowbreak{}MIXING\_\allowbreak{}SUPPORT} & proven & analytic \\
\texttt{\scriptsize G9\_\allowbreak{}HERMITIAN\_\allowbreak{}FOLDED\_\allowbreak{}SUPPORT} & proven & analytic \\
\texttt{\scriptsize G9\_\allowbreak{}LOCAL\_\allowbreak{}SIXTH\_\allowbreak{}NONCANCELLATION} & proven & analytic \\
\texttt{\scriptsize G9\_\allowbreak{}ONE\_\allowbreak{}FACE\_\allowbreak{}SIXTH} & proven & analytic \\
\texttt{\scriptsize G9\_\allowbreak{}SIXTH\_\allowbreak{}ORDER\_\allowbreak{}RUR\_\allowbreak{}ACTIVATION} & disputed & record-backed \\
\texttt{\scriptsize GAUSSIAN\_\allowbreak{}FULL\_\allowbreak{}FAST\_\allowbreak{}GREEN} & proven & analytic \\
\texttt{\scriptsize GAUSSIAN\_\allowbreak{}OS\_\allowbreak{}OBSERVABILITY} & proven & analytic \\
\texttt{\scriptsize GAUSSIAN\_\allowbreak{}PATH\_\allowbreak{}ENDPOINT\_\allowbreak{}BASELINE} & proven & analytic \\
\texttt{\scriptsize GAUSSIAN\_\allowbreak{}QUANTUM\_\allowbreak{}FAST\_\allowbreak{}SOURCES} & proven & analytic \\
\texttt{\scriptsize GAUSSIAN\_\allowbreak{}SCHUR\_\allowbreak{}MEMORY} & proven & analytic \\
\texttt{\scriptsize GROUND\_\allowbreak{}MARGINAL\_\allowbreak{}SCHUR\_\allowbreak{}SCORE} & proven & analytic \\
\texttt{\scriptsize HODGE\_\allowbreak{}FESHBACH\_\allowbreak{}SPLITTING} & proven & analytic \\
\texttt{\scriptsize HODGE\_\allowbreak{}INDUCED\_\allowbreak{}EIGHTH\_\allowbreak{}ORDER\_\allowbreak{}MOMENT} & proven & analytic \\
\texttt{\scriptsize HODGE\_\allowbreak{}LEAKAGE\_\allowbreak{}ZERO\_\allowbreak{}EXTENSION} & proven & analytic \\
\texttt{\scriptsize HODGE\_\allowbreak{}RANK\_\allowbreak{}ONE\_\allowbreak{}SECULAR\_\allowbreak{}CUBIC} & proven & analytic \\
\texttt{\scriptsize INTRINSIC\_\allowbreak{}SCORE\_\allowbreak{}CURRENT} & proven & analytic \\
\texttt{\scriptsize LITERAL\_\allowbreak{}ENDPOINT\_\allowbreak{}SPECTRAL\_\allowbreak{}COMPARISON} & proven & analytic \\
\texttt{\scriptsize LOCAL\_\allowbreak{}CLASS\_\allowbreak{}WICK\_\allowbreak{}COEFFICIENTS} & proven & analytic \\
\texttt{\scriptsize LOCAL\_\allowbreak{}CLASS\_\allowbreak{}WICK\_\allowbreak{}RECURRENCE} & proven & analytic \\
\texttt{\scriptsize LOCAL\_\allowbreak{}CLASS\_\allowbreak{}WICK\_\allowbreak{}SIGNS} & proven & analytic \\
\texttt{\scriptsize LOCAL\_\allowbreak{}GRADIENT\_\allowbreak{}EXCITATION\_\allowbreak{}SUPPORT} & proven & analytic \\
\texttt{\scriptsize MESOSCOPIC\_\allowbreak{}NORMALIZED\_\allowbreak{}SOURCE} & proven & analytic \\
\texttt{\scriptsize MOVING\_\allowbreak{}TIME\_\allowbreak{}MULTICHANNEL\_\allowbreak{}RATE} & proven & analytic \\
\texttt{\scriptsize MOVING\_\allowbreak{}TIME\_\allowbreak{}MULTICHANNEL\_\allowbreak{}SHARPNESS} & proven & analytic \\
\texttt{\scriptsize MOVING\_\allowbreak{}TIME\_\allowbreak{}SPECTRAL\_\allowbreak{}GAP} & proven & analytic \\
\texttt{\scriptsize ODD\_\allowbreak{}ORDER\_\allowbreak{}CENTRE\_\allowbreak{}PARITY} & proven & analytic \\
\texttt{\scriptsize ORDERED\_\allowbreak{}CONTOUR\_\allowbreak{}ACTIVITIES} & proven & analytic \\
\texttt{\scriptsize OS\_\allowbreak{}HISTORY\_\allowbreak{}BLOCK\_\allowbreak{}INTERTWINER} & proven & analytic \\
\texttt{\scriptsize OS\_\allowbreak{}KERNEL\_\allowbreak{}MOVING\_\allowbreak{}TIME\_\allowbreak{}GAP} & proven & analytic \\
\texttt{\scriptsize PREDICTABLE\_\allowbreak{}SOURCE\_\allowbreak{}DRIFT} & proven & analytic \\
\texttt{\scriptsize PROJECTIVE\_\allowbreak{}SOURCE\_\allowbreak{}MARTINGALE} & proven & analytic \\
\texttt{\scriptsize REVERSE\_\allowbreak{}MARTINGALE\_\allowbreak{}BLOCKING} & proven & analytic \\
\texttt{\scriptsize RICCATI\_\allowbreak{}RESIDUAL\_\allowbreak{}CERTIFICATE} & proven & analytic \\
\texttt{\scriptsize ROUGH\_\allowbreak{}GAUGE\_\allowbreak{}COERCIVITY\_\allowbreak{}LOWER\_\allowbreak{}BOUND} & proven & analytic \\
\texttt{\scriptsize RSM\_\allowbreak{}R\_\allowbreak{}CARRIER\_\allowbreak{}SYMBOL\_\allowbreak{}MASTER\_\allowbreak{}THEOREM} & proven & analytic \\
\texttt{\scriptsize SELECTED\_\allowbreak{}INVERSE\_\allowbreak{}WITHOUT\_\allowbreak{}DIAGONAL} & proven & analytic \\
\texttt{\scriptsize SHARP\_\allowbreak{}CUTOFF\_\allowbreak{}SOURCE\_\allowbreak{}GAP\_\allowbreak{}BUDGET} & proven & analytic \\
\texttt{\scriptsize SOURCE\_\allowbreak{}COMPLEMENT\_\allowbreak{}COMPLETE\_\allowbreak{}WINDOW} & proven & analytic \\
\texttt{\scriptsize SOURCE\_\allowbreak{}CONGRUENCE\_\allowbreak{}METRIC\_\allowbreak{}TRANSPORT} & proven & analytic \\
\texttt{\scriptsize SOURCE\_\allowbreak{}CURRENT\_\allowbreak{}BOUNDED\_\allowbreak{}W6} & proven & analytic \\
\texttt{\scriptsize SOURCE\_\allowbreak{}CURRENT\_\allowbreak{}VARIATIONAL\_\allowbreak{}ENERGY} & proven & analytic \\
\texttt{\scriptsize SQUARE\_\allowbreak{}FIRST\_\allowbreak{}CURRENT\_\allowbreak{}BUDGET} & proven & analytic \\
\texttt{\scriptsize SQUARE\_\allowbreak{}FIRST\_\allowbreak{}RESIDUAL\_\allowbreak{}CURRENT} & proven & analytic \\
\texttt{\scriptsize SQUARE\_\allowbreak{}SCALAR\_\allowbreak{}SOURCE\_\allowbreak{}ODD\_\allowbreak{}SCHUR\_\allowbreak{}VANISHING} & proven & analytic \\
\texttt{\scriptsize TETRAHEDRAL\_\allowbreak{}FLUX\_\allowbreak{}PROJECTION\_\allowbreak{}BOUNDARY} & proven & analytic \\
\texttt{\scriptsize TETRAHEDRAL\_\allowbreak{}S4\_\allowbreak{}COMMUTANT\_\allowbreak{}RESOLUTION} & proven & analytic \\
\texttt{\scriptsize TIER\_\allowbreak{}COLLAPSE\_\allowbreak{}ACTUAL\_\allowbreak{}H4\_\allowbreak{}SUPPORT} & proven & analytic \\
\texttt{\scriptsize TIER\_\allowbreak{}COLLAPSE\_\allowbreak{}IS\_\allowbreak{}R\_\allowbreak{}DEGREE} & falsified & analytic \\
\texttt{\scriptsize TRUE\_\allowbreak{}GROUND\_\allowbreak{}CENTER\_\allowbreak{}SCORE\_\allowbreak{}OBSTRUCTION} & proven & analytic \\
\texttt{\scriptsize UNIFORM\_\allowbreak{}MARKED\_\allowbreak{}CLUSTERING} & proven & analytic \\
\texttt{\scriptsize UNIVERSAL\_\allowbreak{}CELLULAR\_\allowbreak{}HODGE\_\allowbreak{}SPECTRUM} & proven & analytic \\
\texttt{\scriptsize W6\_\allowbreak{}ACTUAL\_\allowbreak{}CENTERED\_\allowbreak{}COMPLEMENT\_\allowbreak{}SUPPRESSION} & proven & analytic \\
\texttt{\scriptsize W6\_\allowbreak{}ANGULAR\_\allowbreak{}REFERENCE\_\allowbreak{}SCORE\_\allowbreak{}BUDGET} & proven & analytic \\
\texttt{\scriptsize W6\_\allowbreak{}R10\_\allowbreak{}CONDITIONAL\_\allowbreak{}DERIVATIVE} & proven & analytic \\
\texttt{\scriptsize W6\_\allowbreak{}R10\_\allowbreak{}H4\_\allowbreak{}FROM\_\allowbreak{}H0} & proven & analytic \\
\texttt{\scriptsize W6\_\allowbreak{}R10\_\allowbreak{}PROJECTION\_\allowbreak{}JET\_\allowbreak{}RECURRENCE} & proven & analytic \\
\texttt{\scriptsize W6\_\allowbreak{}UNIFORM\_\allowbreak{}AGMON\_\allowbreak{}COMPLEMENT\_\allowbreak{}GAP} & proven & analytic \\
\texttt{\scriptsize WILSON\_\allowbreak{}ACTIVITY\_\allowbreak{}EXTRACTION} & proven & analytic \\
\texttt{\scriptsize WILSON\_\allowbreak{}ACTUAL\_\allowbreak{}BLOCK\_\allowbreak{}FAST\_\allowbreak{}COMPLEMENT} & proven & analytic \\
\texttt{\scriptsize WILSON\_\allowbreak{}BLOCK\_\allowbreak{}SCORE\_\allowbreak{}OBSTRUCTION} & proven & analytic \\
\texttt{\scriptsize WILSON\_\allowbreak{}CARDINALITY\_\allowbreak{}CHART} & proven & analytic \\
\texttt{\scriptsize WILSON\_\allowbreak{}COEFFICIENT\_\allowbreak{}LOCALITY} & proven & analytic \\
\texttt{\scriptsize WILSON\_\allowbreak{}COMMON\_\allowbreak{}GAUSS\_\allowbreak{}LITERAL\_\allowbreak{}FAST\_\allowbreak{}COMPLEMENT} & proven & analytic \\
\texttt{\scriptsize WILSON\_\allowbreak{}COVARIANT\_\allowbreak{}BLOCK\_\allowbreak{}SOURCES} & proven & analytic \\
\texttt{\scriptsize WILSON\_\allowbreak{}CREATOR\_\allowbreak{}LIMIT} & proven & analytic \\
\texttt{\scriptsize WILSON\_\allowbreak{}CUBIC\_\allowbreak{}GROUND\_\allowbreak{}TRANSFER} & proven & analytic \\
\texttt{\scriptsize WILSON\_\allowbreak{}ENDPOINT\_\allowbreak{}COMPLETE\_\allowbreak{}CLUSTER} & proven & analytic \\
\texttt{\scriptsize WILSON\_\allowbreak{}ENDPOINT\_\allowbreak{}EQUATION} & proven & analytic \\
\texttt{\scriptsize WILSON\_\allowbreak{}FINITE\_\allowbreak{}CELL\_\allowbreak{}PHYSICAL\_\allowbreak{}GAP} & proven & analytic \\
\texttt{\scriptsize WILSON\_\allowbreak{}FIXED\_\allowbreak{}ORDER\_\allowbreak{}CHART} & proven & analytic \\
\texttt{\scriptsize WILSON\_\allowbreak{}FLAT\_\allowbreak{}BACKGROUND\_\allowbreak{}FAST\_\allowbreak{}SOURCES} & proven & analytic \\
\texttt{\scriptsize WILSON\_\allowbreak{}GLOBAL\_\allowbreak{}VERTICAL\_\allowbreak{}BARRIER} & proven & analytic \\
\texttt{\scriptsize WILSON\_\allowbreak{}GROUND\_\allowbreak{}BUNDLE\_\allowbreak{}RELATIVE\_\allowbreak{}FORM} & proven & analytic \\
\texttt{\scriptsize WILSON\_\allowbreak{}HARMONIC\_\allowbreak{}BOUNDARY} & proven & analytic \\
\texttt{\scriptsize WILSON\_\allowbreak{}HARMONIC\_\allowbreak{}CUBIC\_\allowbreak{}OBSTRUCTION} & proven & analytic \\
\texttt{\scriptsize WILSON\_\allowbreak{}INFINITE\_\allowbreak{}PHYSICAL\_\allowbreak{}BAND} & proven & analytic \\
\texttt{\scriptsize WILSON\_\allowbreak{}KINETIC\_\allowbreak{}WINDOW} & proven & analytic \\
\texttt{\scriptsize WILSON\_\allowbreak{}LITERAL\_\allowbreak{}COARSE\_\allowbreak{}FORM} & proven & analytic \\
\texttt{\scriptsize WILSON\_\allowbreak{}LITERAL\_\allowbreak{}FAST\_\allowbreak{}COMPLEMENT} & proven & analytic \\
\texttt{\scriptsize WILSON\_\allowbreak{}LOCALIZED\_\allowbreak{}TRUE\_\allowbreak{}GROUND\_\allowbreak{}SCORE} & proven & analytic \\
\texttt{\scriptsize WILSON\_\allowbreak{}PARENT\_\allowbreak{}GNS} & proven & analytic \\
\texttt{\scriptsize WILSON\_\allowbreak{}PHYSICAL\_\allowbreak{}SOURCE\_\allowbreak{}RANK\_\allowbreak{}REPAIR} & proven & analytic \\
\texttt{\scriptsize WILSON\_\allowbreak{}ROOTED\_\allowbreak{}CONTRACTION} & proven & analytic \\
\texttt{\scriptsize WILSON\_\allowbreak{}SAME\_\allowbreak{}WEIGHT\_\allowbreak{}OBSTRUCTION} & proven & analytic \\
\texttt{\scriptsize WILSON\_\allowbreak{}SECOND\_\allowbreak{}ORDER\_\allowbreak{}CHART} & proven & analytic \\
\texttt{\scriptsize WILSON\_\allowbreak{}STRIP\_\allowbreak{}BO\_\allowbreak{}FIRST\_\allowbreak{}TERM} & proven & analytic \\
\texttt{\scriptsize WILSON\_\allowbreak{}SYMMETRIC\_\allowbreak{}CREATORS} & proven & analytic \\
\texttt{\scriptsize WILSON\_\allowbreak{}THREE\_\allowbreak{}DIMENSIONAL\_\allowbreak{}HARMONIC\_\allowbreak{}BOUNDARY} & proven & analytic \\
\texttt{\scriptsize WILSON\_\allowbreak{}TRANSFER\_\allowbreak{}RESOLVENT} & proven & analytic \\
\texttt{\scriptsize WILSON\_\allowbreak{}TWO\_\allowbreak{}SQUARE\_\allowbreak{}PHYSICAL\_\allowbreak{}SHELLS} & proven & analytic \\
\texttt{\scriptsize WILSON\_\allowbreak{}TWO\_\allowbreak{}STRIP\_\allowbreak{}PHYSICAL\_\allowbreak{}SPLITTING} & proven & analytic \\
\texttt{\scriptsize WILSON\_\allowbreak{}UNIFORM\_\allowbreak{}FINITE\_\allowbreak{}SHELL} & proven & analytic \\
\texttt{\scriptsize WILSON\_\allowbreak{}VACUUM\_\allowbreak{}COMPRESSION} & proven & analytic \\
\texttt{\scriptsize WILSON\_\allowbreak{}VACUUM\_\allowbreak{}SPECTRAL\_\allowbreak{}FLOW} & proven & analytic \\
\texttt{\scriptsize WILSON\_\allowbreak{}VERTICAL\_\allowbreak{}FAST\_\allowbreak{}ENERGY} & proven & analytic \\
\texttt{\scriptsize WILSON\_\allowbreak{}WEIGHTED\_\allowbreak{}ACTIVITIES} & proven & analytic \\
\end{longtable}

